\begin{table}[!htbp]
  \centering
  \small
  \setlength{\tabcolsep}{5pt}
  \begin{tabular}{l cccc}
    \toprule
    \textbf{任务} & Uniform & BA-$k_7$ & Heur-$k_7$ & BA-AutoK (90\%) \\
    \midrule
    NarrativeQA (F1) & \textbf{7.05} & 6.67 & \textbf{6.83} & \textbf{6.80} \\
    HotpotQA (F1) & \textbf{5.57} & \textbf{5.49} & \textbf{5.40} & 5.39 \\
    GovReport (Rouge-L) & \textbf{9.09} & 8.95 & \textbf{9.03} & \textbf{9.27} \\
    \midrule
    \textit{平均} & \textit{7.23} & \textit{7.04} & \textit{7.09} & \textit{7.15} \\
    \bottomrule
  \end{tabular}
  \caption{Qwen2.5-14B 高性能策略簇的任务级分布。\textbf{加粗}表示每个任务的 top-3 策略。Top-3 策略相对 top-1 的最大相对差距为 3.5\%，平均约为 3.0\%，说明 14B 上多种策略位于接近的高性能区间内，当前证据下未形成稳定跨任务最优策略。}
  \label{tab:t6-14b-toptier}
  \par\smallskip\noindent{\footnotesize 相对差距定义为 $(top_1 - top_3) / top_1$。由于差距量级处于 bootstrap CI 半宽范围内，14B 高性能策略簇内部的可区分性需要结合 $n{=}5$ seed CI 解释。在当前统计功效下，应读作不可分辨。}
\end{table}
